\clearpage
\section*{S14. Full-material preconditioner execution}
The observation and injection factors also determine a low-rank change of the material normal matrix. This section follows that change into an inverse used by full-H PCG, then records the checks and costs of its implementation. The final solve continues to use the complete material operator.
The new experiments reuse historical states and preserve their original material fields, damping, data, acquisition, and material parameterization. All material coordinates use the declared volume-weighted physical metric. Complex dielectric increments are packed into real and imaginary parts; because the frozen DDA derivative is complex-linear in the increment, the implementation uses its equivalent complex representation internally. The requested two-column current enrichment is a complex current rank, not a two-dimensional real material repair. Over six illuminations, the update signatures may have many more material columns.

Gate 0 compares a dense inverse, an indefinite Woodbury realization, and rank-revealing QR followed by a small SPD solve. Seven cases include random complex matrices, exact and near rank loss, adjoint and realification checks, and a 27-cell vector DDA. The largest small-DDA inverse residual is approximately \(6.3\times10^{-13}\). A separate adapter check uses the production tangent convention for Gaussian and voxel coordinates, with discrepancies below \(3\times10^{-15}\). These are implementation checks, not continuum validation.

Let \(D=Y_bT_b\). The signature-based identity
\[
H_C-H_0=[J_0^*Y_b,T_b^*]
\begin{bmatrix}0&I\\I&Y_b^*Y_b\end{bmatrix}
[J_0^*Y_b,T_b^*]^*
\]
uses block-diagonal observation factors for different illuminations and a shared material coordinate. The small inverse does not evaluate full \(JZ\) or \(HZ\). Its setup still includes candidate projection/SVD, \(LC\), local injection adjoints, whitening by the baseline inverse, and a small Cholesky factorization. An independently built enlarged Galerkin model checks the resulting Jacobian and inverse. These reconstruction checks are offline validation: their costs appear in job occupation, not in the marginal solver comparison.

Every full-H solve starts at the same old-model step. The stopping target is
\[
\frac{\|Hs+h\|_2}{\|h\|_2}\le10^{-8},
\]
with a maximum of 150 PCG iterations and an explicitly recomputed final residual. Each full normal action entails six primal and six adjoint full-wave right-hand sides. The initial and final residual evaluations are charged. No dense full-H factorization is used by the online solver. Baseline timings include CPU copies that retain the first six paid current/material actions; their instrumentation overhead is not separately attributed. Thus the reported wall ratios concern this instrumented prototype, and small single-replay margins cannot establish production acceleration. Exact small-tangent spectra, reference intervals, and the separate 30-step and 100-step voxel Ritz diagnostic are offline data.

A computable conservative energy-error interval follows from \(H\succeq\lambda I\):
\[
\|s-s_*\|_H\le\|Hs+h\|_2/\sqrt\lambda,
\qquad \Phi(s)-\Phi(s_*)\le\|Hs+h\|_2^2/(2\lambda).
\]
The CSV uses an upper bound for \(\|h\|\) obtained from the stored initial full-H action and initial defect. This can be conservative; it is not labeled an exact reference error.

\section*{S15. Probabilistic certificates and reuse accounting}
The bidirectional tail gives a sufficient spectral condition; evaluating it cheaply requires a separate statistical check. Independent receiver/material actions probe the omitted transfer, with an explicit failure probability. We then account for how checking, construction, and history acquisition affect repeated use of the inverse.
For a fixed real operator \(A\) and independent standard real Gaussian vectors \(g_i\),
\[
\Pr\{\|A\|_2>2\max_{1\le i\le8}\|Ag_i\|_2\}
\le[2\Phi_{\rm N}(1/2)-1]^8=4.62285\times10^{-4}.
\]
The primary campaign uses \(A=(J-J_C)\Lambda^{-1/2}\). Its upper bound also bounds \(\eta_C\), but can be loose. All candidates are frozen before these validation probes. A union bound over the 72 candidate checks is 0.0333; sharing probe vectors across fixed banks does not require cross-bank independence for this union bound. One nonvacuous upper bound is not evidence that the other bounds failed probabilistically: a value above one simply gives no contraction certificate.

A post-primary four-case check replaces $\Lambda$ whitening by the available exact $H_C$ inverse factor in the Jacobian probe operator. For prior-adjoint ranks 2/12, the two objects yield upper bounds 3.765/2.342 and 14.345/8.861. All remain vacuous; the factor and probe costs are recorded. These results exclude an easy resolution merely by changing the whitening.

The reuse check factors \(M_{C,a}=FF^*\) and probes the symmetric operator
\[
A=F^*(H_t-H_{C,a})F.
\]
An upper bound below one yields spectral equivalence; the declared refresh threshold is 0.6. The small factor is constructed from the positive eigenvalues of the QR-SPD block, with no hidden regularization. Eight full-H actions per check are charged. A certificate based on a mean-square probe statistic would not satisfy the stated maximum-probe probability formula and is not used.

Chronological policies are always refresh, never refresh, refresh every two saved states, and independent-probe-triggered refresh. Cached construction timings are recharged at each simulated refresh, so the reported totals are measured phase sums, not uninterrupted end-to-end trajectories. At the first state, no historical adjoint bank exists; the reuse driver records a residual-current fallback. Later adjoint banks use previous-state paid current snapshots.

The frozen-H multiple-RHS experiment uses eight distinct one-percent complex measurement-noise realizations at the same late material state. It represents repeated local uncertainty/data-perturbation solves, not eight nonlinear reconstructions. Full state, baseline inverse, and the corresponding RHS construction are genuinely common. Snapshot acquisition is available from the disclosed preceding full solve; it is not a free cold-start resource. The setup and one probabilistic check are paid before the eight proposed solves. Baseline/proposed order alternates.

The initial rank-two material-recycling control receives already-paid histories from separately timed baseline solves. It is therefore an externally assisted historical-information comparator, not a standalone recycling policy. The rank-12 extension corrects this deployment distinction: after the common initial history, each recycling solve captures its own full-H actions and reuses only those. Both controls use a refreshed material coarse correction followed by full-H PCG. Neither is claimed to exhaust optimized deflated/block/recycling Krylov implementations.

\section*{S16. Structured-path experiment and limits}
The principal path experiment has a one-dimensional chain of 13 wavevectors with all three field components. A uniform lossy background is inverted blockwise through the dyadic Maxwell symbol; a nearest-neighbor modulation defines the residual mode graph. Injection and observation lie in a specified noncentral channel. Increasing the retained radius from zero to three changes the minimum omitted round-trip order from two to eight. The residual coupling norm is explicitly scaled to the reported \(\rho\); this is a controlled finite Fourier model rather than an Ewald periodic simulation of the sharp target.

All four selectors use matched current rank: support-path neighborhoods, centered low frequency, background-propagation block strength, and direct observation strength with deterministic ties. The leading exponent is fitted on the smallest three coupling values; the theoretical statement remains an upper bound. Strong residual coupling outside \(\rho<1\) is retained as an assumption-failure check in the auxiliary periodic stress experiment. An auxiliary worker's control labeled ``radiation'' actually used propagation strength; that row is excluded from receiver-ranking claims.

The finite free-space witness uses the implemented polarizability derivative and total field. One seeded material perturbation produces six illumination-dependent tangent-current responses through the full current equation. The construction projects these responses outside the joint receiver-row and old-current space and retains two independent orthonormal columns. The full material Jacobian is used for evaluation, not to select the witness. This deliberate construction tests existence rather than prevalence or a general online selection rule. The current solve residual stays around \(10^{-15}\), while the feedback-dressed derivative change is percent-level. The uniform-grid changes are a discrete sensitivity check, not independent continuum validation. Supplement S23 specifies the output and derivative norms and their real-material interpretation.

\section*{S17. Reproduction and evidence hierarchy}
The accompanying execution archive provides a frozen protocol, source hashes, raw anchor JSON and paid-current histories, resource logs, failures, and experiment/cost ledgers. The original paired-repair and sharp-reconstruction records remain unchanged. New records are labeled by stage and cannot be pooled with the historical 72/36 comparisons as independent trials. The standalone reproduction guide identifies local CPU tests and the direct-LAN CUDA commands. The code retains FP64, fixed within-solve preconditioning, an explicit SPD failure path, and a serial GPU watchdog. No FFT acceleration or native TSOM optimizer is implemented in this extension; no speed or superiority claim depends on either.
